\def\EditionMode{0}\def\EditionLabel{Diplomatic local review proof}\def\EditionStatus{Local proof}
\documentclass[10pt,oneside]{book}
\usepackage{fontspec}
\setmainfont{Latin Modern Roman}
\setsansfont{Latin Modern Sans}
\usepackage{amsmath,amssymb,mathtools}
\usepackage{unicode-math}
\setmathfont{Latin Modern Math}
\usepackage{xcolor,tikz-cd,lua-ul,marginnote,pdflscape}
\usetikzlibrary{arrows.meta,decorations.pathmorphing,calc,shapes.geometric}
\tikzcdset{ampersand replacement=\&}
\tikzcdset{every label/.append style={text=black}}
\usepackage[a4paper,left=18mm,right=45mm,top=18mm,bottom=22mm,marginparwidth=28mm,marginparsep=7mm]{geometry}
\usepackage[unicode,colorlinks=true,urlcolor=black,linkcolor=black]{hyperref}
\usepackage{bookmark}
\setlength{\parindent}{0pt}
\setlength{\parskip}{4pt plus 1pt minus 1pt}
\setlength{\emergencystretch}{1.5em}
\raggedbottom
\definecolor{PSHand}{RGB}{25,58,132}
\definecolor{PSDelete}{RGB}{146,37,37}
\definecolor{PSTypedOld}{RGB}{95,95,95}
\ifnum\EditionMode=2\colorlet{blue}{black}\else\colorlet{blue}{PSHand}\fi
\RenewDocumentCommand{\textcircled}{m}{\tikz[baseline=(PSCircle.base)]\node[draw,circle,inner sep=.5pt,line width=.4pt](PSCircle){#1};}
\newif\ifPSInApparatus
\newcounter{psnote}
\newcounter{psvariant}
\newcounter{pslinkmark}
\newif\ifPSInLink
\newsavebox{\PSNoteBox}
\newdimen\PSNoteLimit
{\footnotesize\global\PSNoteLimit=3\baselineskip}
\newcommand{\CurrentScanURL}{}
\newcommand{\CurrentKey}{}
\newcommand{\CurrentScan}{}
\newcommand{\pssource}[3]{\psflushmargins\clearpage\gdef\CurrentKey{#1}\gdef\CurrentScan{#2}\gdef\CurrentScanURL{#3}\setcounter{psvariant}{0}\typeout{PS-SOURCE: #1}}
\newcommand{\psdivision}[2]{\pdfbookmark[0]{#1}{#2}\par\noindent{\large\bfseries #1\par}\medskip}
\newcommand{\pslinkmarker}[2]{\stepcounter{pslinkmark}\latelua{texio.write_nl('log','PS-LINK-SHIP|#1|\CurrentKey|\arabic{pslinkmark}|#2')}}
\newcommand{\pslink}[2]{\ifPSInLink\pslinkmarker{#1}{nested}#2\else\begingroup\PSInLinktrue\href{\CurrentScanURL}{\pslinkmarker{#1}{outer}#2}\endgroup\fi}
\newcommand{\scanpage}[1]{\ifnum\EditionMode=2\leavevmode\marginpar{\scriptsize\sffamily\pslink{scanpage}{scan #1}}\else\par\noindent{\small\sffamily\pslink{scanpage}{[scan #1]}}\par\medskip\fi}
\newcommand{\mspage}[1]{\ifnum\EditionMode=2\leavevmode\marginpar{\scriptsize\sffamily p.\,#1}\else{\small\sffamily[p.\,#1]}\fi}
\newcommand{\hand}[1]{{\ifnum\EditionMode=2\color{black}\else\color{PSHand}\fi #1}}
\let\over\relax
\NewDocumentCommand{\over}{O{} m m}{{\textcolor{PSTypedOld}{#2}\ensuremath{^{\!\nearrow}}\if\relax\detokenize{#1}\relax\hand{#3}\else #3\fi}}
\NewDocumentCommand{\inserted}{O{} m}{{\if\relax\detokenize{#1}\relax\hand{\ensuremath{\langle}#2\ensuremath{\rangle}}\else\ensuremath{\langle}#2\ensuremath{\rangle}\fi}}
\newcommand{\deleted}[1]{{\color{PSDelete}\ifmmode\text{[del.: }#1\text{]}\else[\textsf{del.:} #1]\fi}}
\newcommand{\xout}[1]{{\color{PSTypedOld}\ifmmode\text{[x: }#1\text{]}\else[\textsf{x:} #1]\fi}}
\let\PSMathUnderline\underline
\RenewDocumentCommand{\underline}{+m}{\ifmmode\PSMathUnderline{#1}\else\underLine{#1}\fi}
\NewDocumentCommand{\uline}{+m}{\underline{#1}}
\NewDocumentCommand{\uuline}{+m}{\ifmmode\PSMathUnderline{\PSMathUnderline{#1}}\else\underLine{\underLine{#1}}\fi}
\newunderlinetype\PSWaveLine{\leaders\hbox{\lower.7ex\hbox{\tikz[x=1pt,y=1pt,baseline=0pt]{\draw[line width=.3pt] (0,0) .. controls (1,1) and (2,1) .. (3,0) .. controls (4,-1) and (5,-1) .. (6,0);}}}}
\NewDocumentCommand{\uwave}{+m}{{\PSWaveLine #1}}
\newcommand{\unsure}[1]{\pslink{unsure}{[#1\ifmmode\,?\else\,?\fi]}}
\newcommand{\illegible}[1]{\ifmmode\text{\pslink{illegible}{[illeg.\ #1]}}\else\pslink{illegible}{[illeg.\ #1]}\fi}
\NewDocumentCommand{\diagram}{+m}{\ifmmode\vcenter{\hbox{#1}}\else\par\medskip{\centering #1\par}\medskip\psflushmargins\fi}
\newenvironment{diagramblock}{\par\begin{center}}{\end{center}\par}
\newcommand{\PSMathNotes}{}
\makeatletter
\NewDocumentCommand{\psplacednote}{m +m}{%
      \sbox{\PSNoteBox}{\parbox{\marginparwidth}{\footnotesize\raggedright\hand{\textsf{m#1.} #2}}}%
      \ifdim\dimexpr\ht\PSNoteBox+\dp\PSNoteBox\relax<\PSNoteLimit
        \leavevmode\marginpar{\usebox{\PSNoteBox}}%
      \else
        \par\nobreak\smallskip\hrule\nobreak\smallskip
        {\footnotesize\textsf{Marginal note m#1.}\quad\hand{#2}\par}
        \nobreak\smallskip\hrule\smallskip
      \fi}
\newcommand{\psflushmargins}{\ifmmode\else\ifinner\else\ifx\PSMathNotes\@empty\else\let\PSTempNotes\PSMathNotes\global\let\PSMathNotes\@empty\PSTempNotes\fi\fi\fi}
\NewDocumentCommand{\margin}{+m}{%
  \ifPSInApparatus\ifmmode\hand{\text{[marg.: }#1\text{]}}\else\hand{\textsf{[marg.: }#1\textsf{]}}\fi
  \else\stepcounter{psnote}%
    \ifmmode
      \text{\textsuperscript{m\arabic{psnote}}}%
      \begingroup\edef\PSQueue{\endgroup\noexpand\g@addto@macro\noexpand\PSMathNotes{\noexpand\psplacednote{\arabic{psnote}}{\unexpanded{$#1$}}}}\PSQueue
    \else
      \textsuperscript{m\arabic{psnote}}\psplacednote{\arabic{psnote}}{#1}%
    \fi
  \fi}
\AtBeginDocument{\everymath\expandafter{\the\everymath\aftergroup\psflushmargins}\everydisplay\expandafter{\the\everydisplay\aftergroup\psflushmargins}}
\makeatother
\newcommand{\psnumber}[2]{\expandafter\gdef\csname psnum@#1\endcsname{#2}}
\NewDocumentCommand{\var}{+m +m m m}{\ifnum\EditionMode=0 #1\else\ifnum\EditionMode=1\ifmmode\text{\textsuperscript{\tiny\csname psnum@#4\endcsname}}\else\textsuperscript{\tiny\csname psnum@#4\endcsname}\fi\fi #2\fi}
\newcommand{\psapparatusstart}{\par\bigskip\hrule\smallskip\PSInApparatustrue\begingroup\footnotesize\raggedright\textsf{Apparatus for scan \CurrentScan}\par}
\newcommand{\psapparatusend}{\par\endgroup\PSInApparatusfalse}
\NewDocumentCommand{\psentry}{m m m +m +m m}{\par\noindent\textbf{#2. #3.}\quad\textit{Typescript:} #4\quad\textit{Corrected:} \if\relax\detokenize{#5}\relax\textnormal{[omitted]}\else #5\fi\quad #6\par}
\newcommand{\psfrontmatter}{%
\hypersetup{pageanchor=false}\begin{titlepage}\raggedright
{\Large Alexander Grothendieck\par}\vspace{12mm}
{\Huge Pursuing Stacks\par}\medskip{\LARGE 1983\par}\vspace{12mm}
{\Large\EditionLabel\par}\vspace{8mm}
{\large\EditionStatus\par}\vfill
Université de Montpellier, Grothendieck archive, cotes 134-2--134-8.\par
Including the 26 source pages of R. Brown's translation at cote 134-2, sites 115--143.\par
\vspace*{2mm}\end{titlepage}\hypersetup{pageanchor=true}
\chapter*{Editorial statement}
AI-transcribed by Claude Opus 5.5 (Anthropic, effort high); checked against the scans, and this edition built, by GPT-6 Astra (OpenAI, ChatGPT Pro mode); prompts and records by Claude (Anthropic).

The witness consulted and not reproduced is the LaTeX edition of Mateo Carmona with the collaboration of Ulrik Buchholtz (arXiv 2111.01000 v2, 2021), extended from the 2015 edition by ``the scrivener'', CC0 1.0.

The translation in the appendix is R. Brown's, prepared with assistance from L. Breen and J-L. Loday, as his note says; the rights in it belong to its translator or his heirs.

The rights in the original work belong to its author or the author's representatives or heirs. Where we hold rights, in the AI-produced transcription, translation, corrections and typesetting, they are dedicated to the public domain under CC0 1.0. This edition is published for educational, non-profit purposes, and we claim no rights we do not hold. If you hold rights in this work and object to its publication here, please open a GitHub issue and we will take it down.

\section*{How to read this form}
\ifnum\EditionMode=0
The diplomatic form preserves typed and handwritten wording, cancellations, insertions, overwritten readings, underlining, doubts, and original page numbers. Blue marks handwriting; red labels handwritten cancellations; gray labels typewriter cancellations and underlying overwritten text. Angle brackets mark insertions; a superscript arrow separates an overwritten reading from its replacement. Typed insertions and typed overwriting remain plain. The disposition of long marginal notes is identified explicitly.
\else
The corrected text applies the author's insertions, replacements, cancellations, and recorded transpositions. It retains his English and notation, including consistent forms such as ``standart'', ``developped'', ``it's'' for ``its'', and ``A la''. An isolated slip is corrected only when specifically recorded and justified. Mathematical anomalies are preserved and discussed in the guide.
\ifnum\EditionMode=1
Blue continues to identify handwriting. A numbered call identifies each maximal editorial operation, with its full diplomatic and corrected readings and reason in the apparatus following the source page. Internal layers of a wholly replaced or deleted passage are preserved within that entry rather than duplicated as separate entries. These entries are generated from the same variants that govern the running text.
\else
This reading form omits the apparatus, cancelled material, superseded readings, insertion brackets, and color distinctions. It retains the resulting text, marginal notes, underlining, and unresolved doubts. Scan and source-page references appear in the margin.
\fi\fi

Every scan mark and remaining doubt links to its source page in the archive viewer. Doubts are editorial reading marks, not emendations. The source-page breaks identify the archival sequence; a source page may occupy more than one printed page. A partial working edition is explicitly labelled and contains only reviewed source pages.
\clearpage}

\begin{document}
\pssource{ps_134-6_075}{134-6/75}{https://grothendieck.umontpellier.fr/web/viewer.html?file=../134-6.pdf\#page=76}
% Pursuing Stacks (cotes 134-2 to 134-8): archive cote 134-6, site page 75 (PDF page 76), typescript page 393.
% typescript. Draft compared: ps5.tex:5233-5298.
% AI-transcribed by Claude Opus 5.5 (claude-opus-5-5, adaptive thinking, effort high), Message Batches API, batch ps3 (msgbatch_01XdFmuVVBAcgc7EFxdtCqsj), 2026-11-12; image review recorded in records/CHECK_DECISIONS_134_3_4.tsv.
\scanpage{134-6/75}
\mspage{393}
an arbitrary \hand{$\mathcal{A}$} satisfying our assumptions, even when \hand{$\mathcal{A}$} is a topos and
even when it is of the special type A\^{}, as I do not know of any \hand{$\mathcal{M}$} such
that
\[ (*) \qquad \hand{\mathcal{A}^{\mathcal{M}}\simeq\mathcal{A}}_{\mathrm{ab}} \qquad . \]
However, in case \hand{$\mathcal{A}$} = A\^{}, introducing the dual category B=A$^{\mathrm{o}}$ and \hand{$\mathcal{B}$} = B\^{},
we get (see D) below)
\[ (**) \qquad \hand{\mathcal{A}^{\mathcal{M}} \approxeq \mathcal{B}_{\mathcal{M}} \approxeq \mathcal{B}}_{\overset{\hand{\unsure{\illegible{1}}}}{\mathrm{ab}}} \qquad \text{if} \quad \hand{\mathcal{M}} = \text{(Ab)} \quad , \]
hence we do get a canonical isomorphism \xout{provided} (*) provided A=B say
and hence \hand{$\mathcal{A}=\mathcal{B}$} . Let's look at any rate at the simplest case, namely
when A is \xout{the}\inserted[typed]{a} final object in (Cat), hence \hand{$\mathcal{A}$} can be identified with
(Sets), and \hand{$\mathcal{A}$}$_{\mathrm{ab}}$ with (Ab), the identification between the categor\over[typed]{j}{i}es
\xout{\unsure{(Ab) and Hom(}} (Ab) and \uline{Hom}\hand{$_!$}(\xout{A}(Ab),(Ab)) being obtained \inserted[typed]{(we hope !)} by associating
to every object L in (Ab), the functor
\[ F \hand{\longmapsto} L\otimes F \; : \; \text{(Ab)} \hand{\longrightarrow} \text{(Ab)} \qquad . \]
This being so, the external operation (22) of (Ab) on \uline{Hom}\hand{$_!$} \hand{$\approxeq$} (Ab)
can be interpreted (using th\over[typed]{\unsure{a}}{i}s identification) as the interior tensor
product operation in (Ab). On the other hand, the operation $*$ of (23)
is equally interpreted as nothing but the tensor product in (Ab), which
justifies the notation suggesting a tensor product. It would be nice
checking corresponding compatibilities for \deleted{more}\inserted{a} general object\deleted{s} A in (Cat)
satisfying A=A$^{\mathrm{o}}$ , namely a direct sum of one-object cate\over{q}{g}ories A$_{\mathrm{i}}$ defi\xout{\unsure{n}}%
ned each in terms of a \uline{commutative} monoid M$_{\mathrm{i}}$ - I did'nt work it out
myself, sorry !

It should be noted that the relationship between the two \inserted[typed]{exterior} Hom's in
(29) and (31) is essentially the same as between the two exterior tensor-
type operations (26) and (23), Hom(F,K) designating an object in \hand{$\mathcal{N}$} and
\margin{$*$} \uline{Hom}(F,K) an \hand{$\mathcal{N}$}-valued sheaf on \hand{$\mathcal{A}$}, just as F \hand{$*$} L designates an object
\margin{$\circledast$} \xout{of} in \hand{$\mathcal{M}$} and F \ensuremath{\diagram{\begin{tikzpicture}[baseline=(s.base)]\node[inner sep=0pt] (s) {$*$};\draw[color=blue,line width=.5pt] (s.center) circle (.42em);\end{tikzpicture}}} L an \hand{$\mathcal{M}$}-valued cosheaf on \hand{$\mathcal{A}$}; the graphical device of
\uline{underlining} the symbol Hom (used for sheaves) corresponds to the device
\margin{$*$} of \uline{circling} the symbol $*$ (used for cosheaves). With this luxury of
explanations, I hope the notations int\unsure{\over[typed]{o}{r}}\over[typed]{o}{o}duced here are getting through...

\margin{75}

\pssource{ps_134-6_076}{134-6/76}{https://grothendieck.umontpellier.fr/web/viewer.html?file=../134-6.pdf\#page=77}
% Pursuing Stacks (cotes 134-2 to 134-8): archive cote 134-6, site page 76 (PDF page 77), typescript page 394.
% typescript. Draft compared: ps5.tex:5278-5361.
% AI-transcribed by Claude Opus 5.5 (claude-opus-5-5, adaptive thinking, effort high), Message Batches API, batch ps3 (msgbatch_01XdFmuVVBAcgc7EFxdtCqsj), 2026-11-12; image review recorded in records/CHECK_DECISIONS_134_3_4.tsv.
\scanpage{134-6/76}
\mspage{394}

1\over{\unsure{2}}{2}.8. \hand{and 13.8.}

105.    Let's go on with the overall review of abelianization.

D)  \uline{Duality for topoi} \hand{\uline{of the type $A^\wedge$}, \uline{and} \uline{tentative} \uline{generalizations}.}

The main fact it seems which will give rise to duality statements
for topoi of the type $A^\wedge$ is the following, rather familiar one:

\uline{Proposition} 3.   Let A be a small category, \hand{$\mathcal{M}$} a \uline{U}-category \xout{such that}
stable under small direct limits,
\[ (32) \qquad \hand{\varepsilon}_A : A \hand{\hookrightarrow} A^\wedge \]
the canonical inclusion functor.  Then the following functor is an
equivalence of categories:
\[ (3\over[typed]{2}{3}) \qquad F \hand{\longmapsto} F \circ \hand{\varepsilon}_A : \quad \underline{\mathrm{Hom}}_{\hand{!}}(A^\wedge,\hand{\mathcal{M}}) \hand{\longrightarrow} \underline{\mathrm{Hom}}(A,\hand{\mathcal{M}}) \quad . \]

\over[typed]{I}{A}s I am at a loss to give a reference for this standard fact of cate-%
gory theory,  I'll give in guise of a proof the indication that a quasi-
inverse functor for (33) is given by the familiar construction
\[ (34) \qquad i \hand{\longmapsto} i_{\hand{!}} : \quad \underline{\mathrm{Hom}}(A,\hand{\mathcal{M}}) \hand{\longrightarrow} \underline{\mathrm{Hom}}_{\hand{!}}(A^\wedge,\hand{\mathcal{M}}) \quad , \]
where for any functor
\[ i: A \hand{\longrightarrow} \hand{\mathcal{M}} \quad , \]
the functor  \xout{i}\hand{$_!$}
\[ i_{\hand{!}}: A^\wedge \hand{\longrightarrow} \hand{\mathcal{M}} \]
is defined by the formula
\[ (35) \qquad i_{\hand{!}}(F) = \mathop{\underset{\hand{\longrightarrow}}{\mathrm{lim}}}\limits_{a \text{ in } A_{/F}} \; i(a) \qquad . \]
It is readily checked (and we have already used a number of times)  that
this functor admits a right adjoint
\[ (36) \qquad i^{*} : \hand{\mathcal{M}} \hand{\longrightarrow} A^\wedge \quad , \quad i^{*}(x) = ( a \hand{\longmapsto} \mathrm{Hom}(i(a),x)) \quad , \]
and hence $i_{\hand{!}}$ commutes to small direct limits i.e. is in $\underline{\mathrm{Hom}}_{\hand{!}}(\underset{\hand{\unsure{\illegible{1}}}}{A^{\wedge}},\hand{\mathcal{M}})$ ,
hence (34).

\uline{Remark}. If \hand{$\mathcal{M}$} is a pseudocotopos,  then by \inserted[typed]{prop.3 above and} \inserted{by} coro\deleted{\unsure{a|t}}llary 2 of prop.1 (p.385)
the functor
\[ (37) \qquad i \hand{\longmapsto} i^{*} : \quad \underline{\mathrm{Hom}}(A,\hand{\mathcal{M}}) \hand{\longrightarrow} \underline{\mathrm{Hom}}^{\hand{!}}(\hand{\mathcal{M}}, A^\wedge) \]
(which in any case is fully faithful) is equally an equivalence of catego-%
ries.

\margin{76}

\pssource{ps_134-6_077}{134-6/77}{https://grothendieck.umontpellier.fr/web/viewer.html?file=../134-6.pdf\#page=78}
% Pursuing Stacks (cotes 134-2 to 134-8): archive cote 134-6, site page 77 (PDF page 78), typescript page 395.
% typescript. Draft compared: ps5.tex:5343-5420.
% AI-transcribed by Claude Opus 5.5 (claude-opus-5-5, adaptive thinking, effort high), Message Batches API, batch ps3 (msgbatch_01XdFmuVVBAcgc7EFxdtCqsj), 2026-11-12; image review recorded in records/CHECK_DECISIONS_134_3_4.tsv.
\scanpage{134-6/77}
\mspage{395}

The equivalence (33) of prop.3 can be interpreted by saying that the
\xout{functor} canonical functor (32) from A to $A^{\wedge}$ is 2-\uline{universal} , for \inserted[typed]{functors} \xout{maps}
of A into $\underline{U}$-categories stable under small direct limits, taking as ``map\over{s}{s}\unsure{''}
between such categories functors which commute to small direct limits.
Thus, \over{t}{t}he $\underline{U}$-category $A^{\wedge}$\ensuremath{\diagram{\begin{tikzpicture}[baseline=-.3ex,x=1em,y=1ex]\draw[color=blue,line width=.4pt] (0,-.7)--(.55,2.3);\end{tikzpicture}}} may be viewed as \xout{the} ``the'' category deduced from
A by \over[typed]{q}{a}dding arbitrary direct limits (disregarding the direct limits which
may \inserted[typed]{perchance} already exist in A ...).

Taking the duals of the two members of (33), we get \xout{the} an equivalent
statement of prop.3 ; \xout{to the effect}

\uline{Corollary} 1 . Let \hand{$\mathcal{N}$} be a $\underline{U}$-category stable under small inverse limits.
Then the functor
\[ (38) \qquad F \hand{\mapsto} F \circ \hand{\varepsilon}_A{}^{o} : \underline{Hom}^{\hand{!}}(A^{\wedge o}, \hand{\mathcal{N}}) \hand{\to} \underline{Hom}(A^{o}, \hand{\mathcal{N}}) \]
is an equivalence of categories.

In terms of the topos
\[ \hand{\mathcal{A}} = A^{\wedge} , \]
we may interprete the left-hand side of (33) as the category of \hand{$\mathcal{M}$}-valued
\uline{cosheaves} on this topos, which by prop.3 \over[typed]{i}{c}an be interpreted (up to equi-%
valence) as the category of functors from A to \hand{$\mathcal{M}$}. Dually, the left
hand side of (38) can be viewed as the category of \hand{$\mathcal{N}$}-valued \uline{sheaves} on
the topos \hand{$\mathcal{A}$}, which (via the right hand side) can be interpreted up to
equivalence as the category of functors $A^{o} \hand{\to}$ \hand{$\mathcal{N}$}, i.e. as the category of
\hand{$\mathcal{N}$}-valued \uline{presheaves} o\unsure{\over[typed]{f}{n}} A . As A endowed with the coarsest (``chaotic'')
site structure is a generating site for the topos \inserted{$\mathcal{A}$}, the equivalence (38)
\over{may}{may} be viewed as a particular case of the familiar fact, according to
which (up to equivalence) the category of \hand{$\mathcal{N}$}-valued \xout{pre}sheaves on a topos
can be constructed in terms of \hand{$\mathcal{N}$}-valued sheaves on any $\underline{U}$-site defining
this topos. (When the site structure is chaotic, then those sheaves are
just arbitrary \hand{$\mathcal{N}$}-valued presheaves .)

Assume now that the $\underline{U}$-category \hand{$\mathcal{M}$} is stable under \uline{both} types of
small limits (direct and inverse). Then applying (33) for (A,\hand{$\mathcal{M}$}) and
(38) for (B,\hand{$\mathcal{M}$}) where B=$A^{o}$ , we get \over{the}{an} equivalence
\inserted{(39)} \deleted{\xout{Coroll\unsure{ary} \illegible{1}}} \hand{$\delta_A^{\mathcal{M}}$}\hand{:} $\underline{Hom}_{\hand{!}}(A^{\wedge}, \hand{\mathcal{M}}) \hand{\xrightarrow{\approx}} \underline{Hom}^{\hand{!}}(B^{\wedge\hand{o}}\hand{,} \hand{\mathcal{M}})$ ,

i.e. (as announced in yesterday's notes, p.390) we get :

\uline{Corollary} 2. Let A be a small category, B=$A^{o}$ the dual category, \hand{$\mathcal{M}$} any
$\underline{U}$-category stable under small direct and inverse limits. Than the catego-%
ry of \hand{$\mathcal{M}$}-valued \uline{cosheaves} on the topos $A^{\wedge}$ is equivalent \deleted{canonically} to

\margin{77}

\pssource{ps_134-6_078}{134-6/78}{https://grothendieck.umontpellier.fr/web/viewer.html?file=../134-6.pdf\#page=79}
% Pursuing Stacks (cotes 134-2 to 134-8): archive cote 134-6, site page 78 (PDF page 79), typescript page 396.
% typescript. Draft compared: ps5.tex:5403-5493.
% AI-transcribed by Claude Opus 5.5 (claude-opus-5-5, adaptive thinking, effort high), Message Batches API, batch ps3 (msgbatch_01XdFmuVVBAcgc7EFxdtCqsj), 2026-11-12; image review recorded in records/CHECK_DECISIONS_134_3_4.tsv.
\scanpage{134-6/78}
\mspage{396}

the category of \hand{$\mathcal{M}$}-val\unsure{\over[typed]{a}{u}}\over[typed]{e}{e}d \uline{sheaves} on the topos $B^{\wedge}$.

This equivalence, defined up to unique isomorphism, is deduced
from the diagram of canonical equivalences
% Hand-drawn arrows are marked blue; typed objects remain plain.
\diagram{\begin{tikzcd}[column sep=small]
\underline{\mathrm{Hom}}_{\hand{!}}(A^{\wedge},\hand{\mathcal{M}})\hand{\overset{\mathrm{def}}{=}(A^{\wedge})^{\mathcal{M}}} \arrow[dr,color=blue, "\hand{\approx}"] \& \& \hand{(B^{\wedge})_{\mathcal{M}}\overset{\mathrm{def}}{=}}\underline{\mathrm{Hom}}^{\hand{!}}(B^{\wedge\circ}\hand{,}\hand{\mathcal{M}}) \arrow[dl,color=blue, "\hand{\approx}"'] \\
(40) \& \underline{\mathrm{Hom}}(A,\hand{\mathcal{M}}) = \underline{\mathrm{Hom}}(B^{\circ},\hand{\mathcal{M}}) \& ,
\end{tikzcd}}

and depends upon the choice of a quasi-inverse of the second vertical
equivalence in (40) (which choice can be made, via the dual of (35),
via the choice of small inverse limits in \hand{$\mathcal{M}$}).

\uline{R\unsure{\over[typed]{g}{e}}mark.} It is felt that the way we got the equivalence (39) via (40),
the rôle of A and B in it is symmetric. To give a more precise statement,
consider the equivalence deduced from (39) by passing to the dual catego-%
ries of the two members - using the tautological isomorphisms (6) of
page 385, we get an equivalence \deleted{\hand{\unsure{$(\delta_{\illegible{1}}^{\mathcal M})^\circ$}}}
\[ (39')\quad \hand{(\delta^{\mathcal{M}}_{A})^{\circ}}:\ \underline{\mathrm{Hom}}^{\hand{!}}(A^{\wedge\circ},\hand{\mathcal{N}}) \hand{\xrightarrow{\approx}} \underline{\mathrm{Hom}}_{\hand{!}}(B^{\wedge},\hand{\mathcal{N}}) \quad , \quad \text{\over[typed]{w}{W}here } \hand{\mathcal{N}=\mathcal{M}^{\circ}}\ ; \]
th\unsure{\over[typed]{i}{e}}s equivalence is \uline{canonicall\unsure{\illegible{1}}y} \uline{quasi-inverse} \uline{to the equivalence} \unsure{\illegible{1}}\xout{\illegible{1}}
\hand{$\delta^{\mathcal{N}}_{B}$} \uline{in opposite direction}, associated to the pair (B,\hand{$\mathcal{N}=\mathcal{M}^{\circ}$}) instead of
(A,\hand{$\mathcal{M}$}) .

In the rest of this subsection D), we'll elaborate on some particu-%
lar cases of the eq\unsure{\over[typed]{m}{u}}ivalence (39) between cosheaves and sheaves.

\inserted{\uline{Case}}\over{1}{1}\textsuperscript{\uline{o}}). Assume \hand{$\mathcal{M}$} = (Sets), then by prop.1 (p.384) the right hand
side of (39) is canonically equivalent to $B^{\wedge}$ itself, hence we get an
equivalence
\[ (41)\qquad B^{\wedge} \hand{\xrightarrow{\approx}} \underline{\mathrm{Hom}}_{\hand{!}}(A^{\wedge},\ (\mathrm{Sets})) \quad . \]
If we want to keep track of the symmetry aspect described in the remark
above, we may consider \xout{(41) as} the functor (41) as being deduced from a
canonical"pairing" between the categories $A^{\wedge}$ and $B^{\wedge}$
\[ (42)\qquad \hand{\delta_{A}}:\ A^{\wedge}\ \mathrm{x}\ B^{\wedge} \hand{\longrightarrow} (\mathrm{Sets}) \hand{,} \]
which is an object in
\[ \underline{\mathrm{Hom}}_{\hand{!!}}(A^{\wedge},\ B^{\wedge};\ (\mathrm{Sets})) \]
i.e. which commutes to small direct limits in e ach variable. (For this
interpretation, compare the dual statement contained in f\unsure{\over[typed]{b}{o}}rmula (7) of
page 385 - and note that (41), being an equivalence, commutes to
small direct limits i.e. is in a category $\underline{\mathrm{Hom}}_{\hand{!}}(B^{\wedge},\ \underline{\mathrm{Hom}}_{\hand{!}}(\ \ ,\ \ ))$ .)

\margin{78}

\pssource{ps_134-6_079}{134-6/79}{https://grothendieck.umontpellier.fr/web/viewer.html?file=../134-6.pdf\#page=80}
% Pursuing Stacks (cotes 134-2 to 134-8): archive cote 134-6, site page 79 (PDF page 80), typescript page 397.
% typescript. Draft compared: ps5.tex:5476-5549.
% AI-transcribed by Claude Opus 5.5 (claude-opus-5-5, adaptive thinking, effort high), Message Batches API, batch ps3 (msgbatch_01XdFmuVVBAcgc7EFxdtCqsj), 2026-11-12; image review recorded in records/CHECK_DECISIONS_134_3_4.tsv.
\scanpage{134-6/79}
\mspage{397}

This pairing gives rise, in a symmetric way,  to \deleted{a}\inserted{the} functor (41) (which is
an e quivalence)  \uline{and} to a functor
\[ (41\hand{')}\xout{)} \qquad A\hat{\ } \hand{\xrightarrow{\approx}} \underline{Hom}_{\hand{!}}(B\hat{\ }, (Sets)) \]
which (it turns out) is no\over{ne}{ne}  other (up to canonical isomorphism) than
(41) with B replaces by A (and hence A replaced by B ),  and therefore is
equally an e quivalence.  \uline{Thus the pairing }(42)  \uline{between the two topoi}
\uline{A$\hat{\ }$ and B$\hat{\ }$ has the remarkable property that it defines an equivalence of}
\uline{each of these topoi with the category of} (\uline{set-valued}) \uline{cosheaves on the}
\uline{other}.  I do not know of any other example of a pair of topoi related
in such a remarkable way\xout{s},  which we may express by saying that the two
topoi are ``\uline{dual}'' to each other.

We still have to give an explicit expression for the pairing (42),
plus a convenient notation.  I'll write
\[ (43) \qquad \hand{\delta_A}(F,G) \overset{\hand{def}}{=} F \ast G \qquad \text{for  F in A$\hat{\ }$, G in B$\hat{\ }$ } , \]
and I'll use the canonical equivalence (valid for any pair of small
categories A,B  - not necessarily dual to each other -  and any \uline{U}-catego-%
ry  stable under direct limits),  deduced by twofold application of
prop.3 ,  plus the tautological isomorphism similar to (7) p.385  :
\[ (44) \qquad \underline{Hom}_{\hand{!!}}(A^{\hand{\wedge}},B^{\hand{\wedge}}; \hand{\mathcal{M}}) \hand{\xrightarrow{\approx}} \underline{Hom}(AxB, \hand{\mathcal{M}}) \quad , \quad F \hand{\mapsto} Fo(\hand{\varepsilon}_{A} x \hand{\varepsilon}_B) , \]
which shows that (42) is known up to canonical isomorphis\over[typed]{p}{m}  \over[typed]{o}{w}hen we know
it's restriction to the full subcategory AxB = AxA$^{o}$ ,  identifying as
usual an object a of A to it's image in A$\hat{\ }$, and similarly for B.  If
b is an object in A,  we'll denote by b$^{o}$  the same object viewed as an
object of A$^{o}$ = B.  With these conventions (including (43)) we get the
nice formula
\[ (45) \qquad a \ast b^{o} = Hom_A(b,a) \quad ( = Hom_B(a^{o},b^{o}) ) \hand{, \text{for } a,b \text{ in } A,} \]
which has the required symmetry property  -  which,  for general
objects  F in A$\hat{\ }$ and G in B$\hat{\ }$,  can be stated as  a bifunctorial isomor\inserted[typed]{phism}
\[ (46) \qquad F \ast G \hand{\simeq} G \ast F \qquad , \]
where the operation $\ast$  in the first member refers to the pair (A,B),  and
in the second to the pair  (B,A)).

From (45)  we easily deduce the more general formula for F$\ast$G , when
either F or G is in A resp. B ,  namely

\margin{79}

\pssource{ps_134-6_080}{134-6/80}{https://grothendieck.umontpellier.fr/web/viewer.html?file=../134-6.pdf\#page=81}
% Pursuing Stacks (cotes 134-2 to 134-8): archive cote 134-6, site page 80 (PDF page 81), typescript page 398.
% typescript. Draft compared: ps5.tex:5529-5604.
% AI-transcribed by Claude Opus 5.5 (claude-opus-5-5, adaptive thinking, effort high), Message Batches API, batch ps3 (msgbatch_01XdFmuVVBAcgc7EFxdtCqsj), 2026-11-12; image review recorded in records/CHECK_DECISIONS_134_3_4.tsv.
\scanpage{134-6/80}
\mspage{398}

(47) \[ a * G \hand{\simeq} G(a) \quad , \quad F \underset{\hand{\unsure{\illegible{1}}}}{*} b^{o} \underset{\hand{\unsure{\illegible{1}}}}{\hand{\simeq}} F(b) \quad . \]

\uline{Remarks}.  \unsure{\over[typed]{We}{As}} we are mainly interested \inserted[typed]{here} in abelianization\xout{,} and  (commutati-%
ve) homology and cohomology,  and hence in \xout{commutative} sheaves and
cosheaves with values in additive (even abelian) categories,  we are not
going to use for the time being the relationship between $A^{\wedge}$\xout{and} and $B^{\wedge}$
just touched upon.  We could elaborate a great deal more on it,  for
instance introducing a canonical pairing \inserted[typed]{(more accurately, a bi-sheaf)} with opposite variance to (42)
(48) \[ A^{\wedge o} \mathrm{x} B^{\wedge o} \hand{\longrightarrow} (\mathrm{Sets}) \]
(or what amounts to the same,  canonical functors adjoint to each other
\[ A^{\wedge o} \hand{\to} B^{\wedge} \quad , \quad B^{\wedge o} \hand{\to} A^{\wedge} \quad ) \ , \]
deduced (via the \xout{isom} equivalence dual to (44)
\[ \underline{\mathrm{Hom}}^{\hand{!!}}(A^{\wedge o},B^{\wedge o};\hand{\mathcal{N}}) \overset{\unsure{\approx}}{\hand{\longrightarrow}} \underline{\mathrm{Hom}}(A^{o}\mathrm{x} B^{o},\hand{\mathcal{N}}) \quad ) \]
from the co-pairing
\[ A^{o} \mathrm{x} B^{o} = B \mathrm{x} A \hand{\longrightarrow} (\mathrm{Sets}) \]
given by
\[ (b^{o},a) \hand{\longrightarrow} \mathrm{Hom}(b,a) \qquad (\text{for } a,b \ \text{ in } A\ ) \ . \]
\margin{other description (48) via $\times$ cf 397}
The two pairings (42) and (48) can be given a common interpretation as
Hom  sets in a suitable category
\[ \hand{\mathcal{E}} = \hand{\mathcal{E}}(A) \quad , \]
which is the union of the two full subcategories  $A^{\wedge}$ and $B^{\wedge o}$  (which may
be interpreted as deduced from A by adjoining respectively small direct
and small inverse limits to it) intersecting in the common subcategory A,
the Hom sets in the two directions between an object F of $A^{\wedge}$ and an object
$G^{o}$ of $B^{\wedge o}$ (corresponding to an object G of $B^{\wedge}$ )  being given \xout{by the} respec-%
\xout{formulas} tively by the pairings (48) and (42).  The full relationship \margin{$\times$ cf p. 40\unsure{4|9}-405}
between these pairings is most conveniently expressed,  it seems,  by the
composition law of \xout{homo} maps in $\hand{\mathcal{E}}(A)$,  and associativity for this law.
T\over[typed]{h}{h}e symmetry of the situation with respect to the pair (A,B)  is expressed
by the canonical isomorphism of categories
\[ \hand{\mathcal{E}}(A)^{o} \hand{\simeq} \hand{\mathcal{E}}(B) \qquad (\ \text{where } B = A^{o}\ ) \quad . \]

\inserted{\uline{Case}} 2\textsuperscript{\uline{o}}).  Of direct relevance for the abelianization story is the particu-%
lar case of the equivalence (39),  obtained by taking
\[ \hand{\mathcal{M}} = (\mathrm{Ab}) \quad . \]
Using formula(12\over{)}{)}(page 387) for the pair $(A^{\wedge},(\mathrm{Ab}))$  and formula (14)
for $B^{\wedge}$ in guise of $\hand{\mathcal{A}}$,  we get  the canonical equivalence

\margin{80}

\end{document}
